\begin{frame}{``정답''은 언제나 사람이 정한다: 라벨의 책임 I}
  $y^{(37)}$ 이 ``그 집의 \emph{실제} 거래가''라고 썼는데, ``실제''를
  오래 바라보자 ---
  \begin{itemize}
    \item $y^{(37)}$ 은 ``세상에서 정말 그 집의 가치''가 아니라,
      \kb{그 거래가 기록된 숫자}. 특수 사유(가족 간 거래, 급매)로
      비정상적으로 싸게 거래됐다면 ``가치''는 달라진다.
    \item 즉 라벨은 ``세상의 진실''이 아니라
      \textbf{``그 시점, 그 상황에서 사람이 기록한 값''}.
  \end{itemize}
  \vskip0.5em
  \begin{block}{실무 결과 1: 라벨의 정의는 팀 전체가 합의해야 한다}
    ``이탈 = 30일 이상 미접속''으로 정한 팀과 ``이탈 = 구독 해지''로
    정한 팀이 같은 데이터셋에 라벨을 매기면, \textbf{같은 고객에게 두
    라벨이 동시에} 붙을 수 있다. 정의가 일관성이 없으면 아무리 좋은
    $h$ 도 그 모순을 학습할 수 없다.
  \end{block}
\end{frame}
